\paragraph{Relation to Select-and-Project:} Concurrent work by B\u{a}rb\u{a}lau et al.\ \cite{ref_sp} (S\&P Top-K) also argues for encoder-centric SAE interventions over the traditional decoder-centric use of SAEs, but the mechanism and evaluation differ from ours in several respects. S\&P Top-K selects encoder features aligned with a sensitive attribute and then applies a \emph{hard orthogonal projection} in the model's native embedding space, $V = I_d - \alpha A(A^\top A)^{-1}A^\top$, which removes the entire component of the embedding along the selected attribute subspace with a single global scale $\alpha \in [0,1]$; the intervention is a fixed geometric projection applied once, independent of how strongly any feature actually fires on a given input. Our method instead performs an \emph{activation-weighted subtraction} at a late transformer layer, $h'_t = h_t - \alpha\sum_{k\in\mathcal{F}} m_k (z_t)_k E_k$, where the magnitude removed for each feature is scaled per-token by that feature's own activation $(z_t)_k$ and gated by a threshold $\tau$: tokens on which a bias feature does not fire receive no edit, rather than being projected regardless of relevance. S\&P Top-K is evaluated on vision-language fairness benchmarks (CelebA, FairFace) and on LLM aggressiveness/sycophancy via LLM-as-judge scoring for Llama-3-8B-Instruct; it does not evaluate on text-only stereotype benchmarks such as CrowS-Pairs or StereoSet, nor on Phi-3 or Gemma. Their ablation contrasts encoder-based selection with projection against decoder-based \emph{masked reconstruction} -- two mechanisms that differ in more than just the direction source -- whereas we hold the subtraction mechanism fixed and swap only encoder rows $E_k$ for decoder columns $D_{:k}$, isolating the encoder-vs-decoder direction question as a controlled ablation. We view the two approaches as complementary evidence for the same broader claim -- that encoder feature directions are a more faithful handle on bias-aligned structure than decoder reconstruction directions -- obtained through different intervention mechanisms (hard projection vs.\ activation-weighted subtraction) and disjoint evaluation domains (vision-language fairness and LLM-judged generation vs.\ token-likelihood stereotype benchmarks).
